\chapter*{Introduction}
\addcontentsline{toc}{chapter}{Introduction}

Le but de ce mémoire est de présenter les travaux réalisés dans le cadre de mon stage de fin d'études, effectué au sein du laboratoire de recherche LATTICE\footnote{Le Lattice est un laboratoire du CNRS (UMR8094), rattaché à l’Ecole Normale Supérieure (IDEX PSL : Paris Sciences et Lettres) et à l’Université Sorbonne Nouvelle-Paris 3. Il regroupe des chercheurs CNRS, des enseignants-chercheurs de l’Université Paris Sorbonne Nouvelle et d’autres universités, ainsi que de nombreux autres personnels permanents ou non-permanents.
Les recherches du Lattice sont consacrées à la linguistique et au traitement automatique des langues.Dans le domaine linguistique proprement dit, nos recherches portent, d’une part, sur l’interface entre grammaire et lexique (dans le courant des grammaires cognitives), et d’autre part, sur le discours (en particulier la structuration du discours et la coréférence).
Le Lattice accorde également une place majeure à l’étude du changement linguistique, qu’il s’agisse de l’évolution du français – en particulier dans les domaines syntaxique et sémantique – ou des tendances plus générales à l’œuvre dans l’évolution des langues.
Dans le domaine du TAL, le Lattice explore les techniques neuronales depuis plusieurs années, et a récemment mis l’accent sur l’analyse syntaxique. Les travaux en linguistique et en TAL sont étroitement corrélés.
Enfin notre laboratoire est fortement impliqué dans le domaine des Humanités Numériques et contribue à la construction de ressources enrichies. https://www.lattice.cnrs.fr/ }. Au cours de ce stage, j’ai travaillé sur de nombreux aspects de \textit{TAL} (Traitement Automatique des Langues). Mon objectif principal était de déterminer comment était redigée l'ecriture naturaliste. Ce stage a été encadré par le professeur \textit{Dominque Legallois} ainsi que par d'autre chercheurs du laboratoire à qui j'ai pu solicité leur aide. 

Ce projet porte plus précisement sur Origine, circularité et transformation des motifs naturalistes. Le but de ce projet est d'analyser les motifs naturalistes dans les textes littéraires et de comprendre comment ils sont utilisés pour créer des effets stylistiques et narratifs. Ce travail à commencé avec la redaction d'un article scientifique sur le sujet par le professeur \textit{Dominique Legallois} 


